\documentclass[10pt]{article}

\usepackage[utf8]{inputenc}
\usepackage[spanish,mexico]{babel}
\usepackage{amsmath,amsfonts,mathabx,amssymb,dsfont}
\usepackage{bbm}
\usepackage{cancel}
\usepackage[colorlinks = true, linkcolor = magenta]{hyperref}
\usepackage[dvipsnames]{xcolor}
\usepackage{wasysym}
\usepackage{graphicx}
\usepackage{listings}
\usepackage{float}

%Pendientes:
%---Buscar implementaciones de modelos sobre cuestiones de derecho laboral
%---Comenzar con encuestas sobre ley laboral
%---Definir bien arquitectura,  ¿BERT-Multilingue?


\author{Pedro D\'{i}az L\'{o}pez}
\title{Fase 1, Introducci\'{o}n a los LLM's}

\lstdefinestyle{mystyle}{
    backgroundcolor= \color{white},   
    commentstyle= \color{purple},
    keywordstyle= \color{magenta},
    numberstyle=\tiny\color{gray},
    stringstyle=\color{blue},
    basicstyle=\ttfamily\footnotesize,
    breakatwhitespace=false,         
    breaklines=true,                 
    captionpos=b,                    
    keepspaces=true,                 
    numbers=left,                    
    numbersep=5pt,                  
    showspaces=false,                
    showstringspaces=false,
    showtabs=false,                  
    tabsize=2
}

\lstset{style=mystyle}


\begin{document}

\maketitle

\subsection{Tarea lingüistica}
\normalsize{Planteamos} para nuestro proyecto un LLM, experto en Derecho Laboral, al que el usuario
pueda realizarle consultas respecto a cuestiones que envuelven las relaciones laborales.
Algunos ejemplos de los posibles cuestionamientos son:
\begin{itemize}
\item Mi jefe me dijo que nos daba el d\'{i}a y a\'{u}n as\'{i} me lo descont\'{o}, ¿Es
eso legal?
\item Me contrataron como t\'{e}cnico de mantenimiento, pero me cobran por el uso de la herramienta
¿Hay algo que pueda hacer?
\end{itemize}

\subsection{Corpus}
Para nuestro proposito, nuestra fuente primaria de informaci\'{o}n ser\'{a} la Ley Federal del
Treabajo, texto disponible en linea en formatos PDF y .doc, además de apoyarnos en otras leyes
tales como el C\'{o}digo de Comercio, seg\'{u}n sea necesario. Naturalmente, \'{e}stos textos
pertenecen al dominio legal 
\subsection{Arquitectura}
Como arquitecura base, proponemos usar un encoder-only, pues, dado el proposito de nuestro LLM,
y como cualquier persona lo sabe, en derecho, el contexto y la forma constituyen el fondo.
\subsection{Baseline}
Como Baseline, esperamos que el modelo se comporte como cualquier persona ajena al ambito
legal al tratar de leer una Ley, que sea capaz de relacionar los cuestionamientos con los
ariculos relacionados, pero no necesariamente será capaz de dar las respuestas adecuadas a los mismos.
\subsection{Motivaci\'{o}n}
Considero que este proyecto es importante por el simple hecho de que el mexicano promedio desconoce
tanto sus derechos, como sus obligaciones en materia laboral, dando paso a un abuso constante
y sist\'{e}mico en las relaciones laborales en las que se ve involucrado. Tener acceso facil
a un experto en materia laboral, podr\'{i}a ayudar a cambiar \'{e}sta din\'{a}mica.	
\end{document}